\section{SQL Analysis and Results}
\label{sec:sql}

The query portfolio contains thirteen analytic queries (six basic and seven
advanced) plus one performance-tuned query and one persisted view. Together
they cover five SQL categories called out by the project rubric: subqueries,
common table expressions (CTEs), window functions, multi-table joins
($\geq 3$ tables), and indexing / \texttt{EXPLAIN} performance work. Every
query is committed in the project repo under \texttt{sql/queries/} with the
naming convention \texttt{<owner>\_<short\_name>.sql}.

For each query we present the domain question, the SQL, a result snapshot,
and a one-paragraph insight. Result snapshots in this report are
PostgreSQL output exported from \texttt{psql} or rendered from the
Streamlit dashboard.

\subsection{Methodology}
All queries run against the loaded database documented in
Section~\ref{sec:dbsetup}. Each query is reproducible: a reader who has
followed the load instructions can paste any \texttt{SELECT} below into
their \texttt{psql} session and obtain the same result. Where helpful we
note the technique label (CTE, window function, etc.) so the rubric
mapping is unambiguous.

% =========================================================================
\subsection{Basic Queries}
% =========================================================================

\subsubsection{B1 (River): Labeled lines per source host}
\textbf{Question.} Which hosts contribute the most labeled attack
evidence?

\begin{lstlisting}[language=SQL]
SELECT source_host,
       COUNT(*) AS labeled_line_count
FROM labeled_line
GROUP BY source_host
ORDER BY labeled_line_count DESC;
\end{lstlisting}

% TODO: insert result snapshot at figures/q_b1_result.png
\textbf{Insight.} % TBD: river to add a 2-3 sentence interpretation.
The dominant host (by labeled-line count) frames the rest of the analysis.

% -------------------------------------------------------------------------
\subsubsection{B2 (River): Attack labels per phase}
\textbf{Question.} How many distinct labels does each attack phase carry?
\begin{lstlisting}[language=SQL]
SELECT ap.phase_name,
       COUNT(al.label_id) AS label_count
FROM attack_phase ap
FULL OUTER JOIN attack_label al ON al.phase_id = ap.phase_id
GROUP BY ap.phase_id, ap.phase_name
ORDER BY label_count DESC;
\end{lstlisting}
% TODO: insert result snapshot at figures/q_b2_result.png
\textbf{Insight.} % TBD river: which phases have the richest taxonomy?

% -------------------------------------------------------------------------
\subsubsection{B3 (River): Audit-event count per host}
\textbf{Question.} Which hosts produce \texttt{auditd} traces, and how
many?

\begin{lstlisting}[language=SQL]
SELECT h.hostname,
       h.host_key,
       COUNT(ae.event_id) AS event_count
FROM host h
INNER JOIN audit_event ae ON ae.host_id = h.host_id
GROUP BY h.host_id, h.hostname, h.host_key
ORDER BY event_count DESC;
\end{lstlisting}
% TODO: insert result snapshot at figures/q_b3_result.png
\textbf{Insight.} % TBD river: confirms only two hosts have audit logs in
% the russellmitchell slice.

% -------------------------------------------------------------------------
\subsubsection{B4 (Ishaan): Labeled lines per attack phase}
\textbf{Question.} How heavily is each phase represented in the labeled
data?
\begin{lstlisting}[language=SQL]
SELECT ap.phase_name, COUNT(*) AS labeled_lines
FROM attack_phase ap
JOIN attack_label al  ON al.phase_id  = ap.phase_id
JOIN labeled_line_label lll ON lll.label_id = al.label_id
GROUP BY ap.phase_name
ORDER BY labeled_lines DESC;
\end{lstlisting}
% TODO: insert result snapshot at figures/q_b4_result.png
\textbf{Insight.} % TBD ishaan: exfiltration dominates by line count.

% -------------------------------------------------------------------------
\subsubsection{B5 (Ishaan): Top-10 detection rules}
\textbf{Question.} Which rules contribute the most labels?
\begin{lstlisting}[language=SQL]
SELECT rule_name, COUNT(*) AS times_used
FROM labeled_line_rule
GROUP BY rule_name
ORDER BY times_used DESC
LIMIT 10;
\end{lstlisting}
% TODO: insert result snapshot at figures/q_b5_result.png
\textbf{Insight.} % TBD ishaan: dnsteal.domain.match dominates.

% -------------------------------------------------------------------------
\subsubsection{B6 (Naman): Service activity baseline}
\textbf{Question.} What is the baseline service activity on monitored
hosts? Establishing the baseline lets us see why one stop/start pair
later stands out as anomalous.
\begin{lstlisting}[language=SQL]
SELECT h.host_key,
       am.unit AS service_name,
       SUM(CASE WHEN ae.type = 'SERVICE_START' THEN 1 ELSE 0 END) AS starts,
       SUM(CASE WHEN ae.type = 'SERVICE_STOP'  THEN 1 ELSE 0 END) AS stops
FROM audit_event ae
JOIN audit_service_event ase ON ase.event_id = ae.event_id
JOIN audit_message am        ON am.event_id  = ae.event_id
JOIN host h                  ON h.host_id    = ae.host_id
GROUP BY h.host_key, am.unit
ORDER BY h.host_key,
         SUM(CASE WHEN ae.type = 'SERVICE_START' THEN 1 ELSE 0 END)
       + SUM(CASE WHEN ae.type = 'SERVICE_STOP'  THEN 1 ELSE 0 END) DESC;
\end{lstlisting}
% TODO: insert result snapshot at figures/q_b6_result.png
\textbf{Insight.} The query returns 33 (host, service) rows.
\texttt{phpsessionclean} dominates \texttt{intranet\_server} with 384
events. The \texttt{put} service on \texttt{internal\_share}, with a
single start and a single stop amid hundreds of routine cron-driven
service events, is a strong candidate for further investigation. This
exact pair is the exfiltration service confirmed by the attack labels.

% =========================================================================
\subsection{Advanced Queries}
% =========================================================================

\subsubsection{A1 (Ishaan): Distinct event types touching labeled hosts (Subquery)}
\textbf{Category.} Subquery (\texttt{EXISTS}).\\
\textbf{Question.} Which audit event types appear on hosts that have at
least one labeled log line?
\begin{lstlisting}[language=SQL]
SELECT DISTINCT e.type
FROM audit_event e
WHERE EXISTS (
  SELECT 1
  FROM host h
  JOIN labeled_line ll ON ll.source_host = h.host_key
  WHERE h.host_id = e.host_id
);
\end{lstlisting}
% TODO: insert result snapshot at figures/q_a1_result.png
\textbf{Insight.} % TBD ishaan: confirms which event types coexist with
% the attack-labeled lines.

% -------------------------------------------------------------------------
\subsubsection{A2 (River): Cross-host distribution of labeled lines (multi-CTE + CROSS JOIN)}
\textbf{Category.} CTE.\\
\textbf{Question.} What share of labeled lines comes from each
(host, log) pair?

\begin{lstlisting}[language=SQL]
WITH per_source AS (
    SELECT source_host, source_log, COUNT(*) AS line_count
    FROM labeled_line
    GROUP BY source_host, source_log
),
total AS (
    SELECT SUM(line_count) AS total_count FROM per_source
)
SELECT ps.source_host,
       ps.source_log,
       ps.line_count,
       ROUND(100.0 * ps.line_count / t.total_count, 2)
         AS pct_of_all_labeled
FROM per_source ps
CROSS JOIN total t
ORDER BY ps.line_count DESC;
\end{lstlisting}
% TODO: insert result snapshot at figures/q_a2_result.png
\textbf{Insight.} % TBD river: shows the long-tail concentration.

% -------------------------------------------------------------------------
\subsubsection{A3 (Ishaan): Busiest day per host (Window function)}
\textbf{Category.} Window function (\texttt{RANK}~\texttt{OVER PARTITION}).\\
\textbf{Question.} On what day did each host generate the most events?

\begin{lstlisting}[language=SQL]
SELECT h.host_key,
       DATE(e.timestamp) AS event_date,
       COUNT(*) AS events_that_day,
       RANK() OVER (
           PARTITION BY h.host_key
           ORDER BY COUNT(*) DESC
       ) AS day_rank_for_host
FROM audit_event e
JOIN host h ON h.host_id = e.host_id
GROUP BY h.host_key, DATE(e.timestamp)
ORDER BY h.host_key, day_rank_for_host, event_date;
\end{lstlisting}
% TODO: insert result snapshot at figures/q_a3_result.png
\textbf{Insight.} % TBD ishaan: daily concentration on the attack day.

% -------------------------------------------------------------------------
\subsubsection{A4 (Naman): Privilege-escalation reconstruction (CTE + window + 6-table join)}
\textbf{Category.} CTE + Window + Multi-table join ($\geq 3$ tables).\\
\textbf{Question.} Can we reconstruct the exact privilege-escalation
attack sequence on \texttt{intranet\_server}?

\begin{lstlisting}[language=SQL]
WITH attack_timeline AS (
    SELECT ae.event_id,
           ae.timestamp,
           ae.type AS audit_type,
           am.op   AS pam_operation,
           am.acct AS target_account,
           am.exe  AS executable,
           am.terminal,
           al.label_name,
           ap.phase_name
    FROM audit_event ae
    JOIN host h          ON h.host_id = ae.host_id
    JOIN labeled_line ll ON ll.source_host = h.host_key
                        AND ll.source_log  = 'audit.log'
                        AND ll.line_number = ae.line_number
    JOIN labeled_line_label lll ON lll.labeled_line_id = ll.labeled_line_id
    JOIN attack_label al ON al.label_id = lll.label_id
    JOIN attack_phase ap ON ap.phase_id = al.phase_id
    LEFT JOIN audit_message am ON am.event_id = ae.event_id
    WHERE h.host_key = 'intranet_server'
)
SELECT event_id, timestamp, audit_type, pam_operation,
       target_account, executable, terminal,
       string_agg(label_name, ', ' ORDER BY label_name) AS labels
FROM attack_timeline
GROUP BY event_id, timestamp, audit_type, pam_operation,
         target_account, executable, terminal
ORDER BY timestamp, event_id;
\end{lstlisting}
% TODO: insert result snapshot at figures/q_a4_result.png
\textbf{Insight.} The query returns nine rows that decompose into two
distinct attack bursts: at 04:37:40 UTC, a \texttt{su} from
\texttt{www-data} to \texttt{jhall} via \texttt{/bin/su} (four events);
then at 04:38:06 UTC, \texttt{sudo cat /etc/shadow} via
\texttt{/usr/bin/sudo} (five events), 25.6 seconds later. The
reconstruction is only possible because audit events, host inventory,
and attack labels live in the same database with cross-domain foreign
keys.

% -------------------------------------------------------------------------
\subsubsection{A5 (Naman): Detection-rule effectiveness (CTE + window)}
\textbf{Category.} CTE + Window function (\texttt{DENSE\_RANK}).\\
\textbf{Question.} Which detection rules are most effective at
identifying attacker activity?

\begin{lstlisting}[language=SQL]
WITH rule_stats AS (
    SELECT llr.rule_name,
           ap.phase_name,
           COUNT(DISTINCT llr.labeled_line_id) AS lines_triggered,
           COUNT(DISTINCT llr.label_id)        AS labels_triggered
    FROM labeled_line_rule llr
    JOIN attack_label al ON al.label_id = llr.label_id
    JOIN attack_phase ap ON ap.phase_id = al.phase_id
    GROUP BY llr.rule_name, ap.phase_name
)
SELECT rule_name,
       phase_name,
       lines_triggered,
       labels_triggered,
       DENSE_RANK() OVER (ORDER BY lines_triggered DESC)
         AS effectiveness_rank
FROM rule_stats
ORDER BY lines_triggered DESC
LIMIT 15;
\end{lstlisting}
% TODO: insert result snapshot at figures/q_a5_result.png
\textbf{Insight.} \texttt{dnsteal.domain.match} is rank 1 with $\sim 53$K
lines, which corresponds to the DNS-tunneling exfiltration phase.
Audit-based rules rank much lower by raw volume but they are the only
ones that catch the privilege-escalation phase. Volume-based ranking is
not a substitute for kill-chain coverage.

% -------------------------------------------------------------------------
\subsubsection{A6 (Naman): Lateral-movement timeline (CTE + 6-table join)}
\textbf{Category.} CTE + Multi-table join ($\geq 3$ tables).\\
\textbf{Question.} Can we trace the attacker's movement \emph{across}
hosts?

\begin{lstlisting}[language=SQL]
WITH labeled_audit AS (
    SELECT h.host_key,
           ae.timestamp,
           ae.type AS audit_type,
           am.acct AS target_account,
           am.exe  AS executable,
           am.unit AS service_unit,
           ap.phase_name,
           string_agg(al.label_name, ', ' ORDER BY al.label_name) AS labels
    FROM audit_event ae
    JOIN host h          ON h.host_id = ae.host_id
    JOIN labeled_line ll ON ll.source_host = h.host_key
                        AND ll.source_log  = 'audit.log'
                        AND ll.line_number = ae.line_number
    JOIN labeled_line_label lll ON lll.labeled_line_id = ll.labeled_line_id
    JOIN attack_label al ON al.label_id = lll.label_id
    JOIN attack_phase ap ON ap.phase_id = al.phase_id
    LEFT JOIN audit_message am ON am.event_id = ae.event_id
    GROUP BY h.host_key, ae.timestamp, ae.type,
             am.acct, am.exe, am.unit, ap.phase_name
)
SELECT host_key, timestamp, audit_type, phase_name, labels,
       COALESCE(target_account, service_unit, '(n/a)') AS target_detail,
       executable
FROM labeled_audit
ORDER BY timestamp, host_key;
\end{lstlisting}
% TODO: insert result snapshot at figures/q_a6_result.png
\textbf{Insight.} The query returns 11 rows across two hosts. At
04:37--04:38 UTC, \texttt{intranet\_server} shows
privilege-escalation activity. At 13:50 UTC, \texttt{internal\_share}
shows exfiltration activity (the \texttt{put} service from query B6).
The two bursts are roughly nine hours apart, on two different machines:
the database makes the lateral movement legible in a single ordering.

% -------------------------------------------------------------------------
\subsubsection{A7 (Naman): Burst detection via timing (CTE + LAG window function)}
\textbf{Category.} CTE + Window function (\texttt{LAG},
\texttt{ROW\_NUMBER}).\\
\textbf{Question.} Can we identify distinct attack bursts from purely
temporal patterns, without relying on any label?

\begin{lstlisting}[language=SQL]
WITH labeled_audit_events AS (
    SELECT ae.event_id, ae.timestamp, ae.type,
           ROW_NUMBER() OVER (ORDER BY ae.timestamp,
                                       ae.line_number) AS seq,
           LAG(ae.timestamp) OVER (ORDER BY ae.timestamp,
                                            ae.line_number) AS prev_timestamp
    FROM audit_event ae
    JOIN host h          ON h.host_id = ae.host_id
    JOIN labeled_line ll ON ll.source_host = h.host_key
                        AND ll.source_log  = 'audit.log'
                        AND ll.line_number = ae.line_number
    WHERE h.host_key = 'intranet_server'
)
SELECT seq, timestamp, type,
       EXTRACT(EPOCH FROM timestamp - prev_timestamp)
         AS seconds_since_prev,
       CASE
           WHEN prev_timestamp IS NULL THEN 'FIRST EVENT'
           WHEN EXTRACT(EPOCH FROM timestamp - prev_timestamp) > 10
                THEN 'NEW BURST'
           ELSE 'continuation'
       END AS burst_indicator
FROM labeled_audit_events
ORDER BY seq;
\end{lstlisting}
% TODO: insert result snapshot at figures/q_a7_result.png
\textbf{Insight.} Two distinct bursts emerge purely from the gap between
consecutive timestamps. Burst 1 (events 1-4): all \texttt{su} steps
within 24 milliseconds of each other. Burst 2 (events 5-9): the sudo
chain, starting 25.6 seconds after the previous event. Without ever
querying the label table, the timing alone identifies two attack
phases.

% =========================================================================
\subsection{Performance: Indexing and EXPLAIN}
\label{sec:explain}
% =========================================================================

The most expensive query in our portfolio is A6 (lateral-movement
timeline): six joins, including the cross-domain join from
\texttt{audit\_event.line\_number} to
\texttt{labeled\_line.line\_number} via the natural-key triple
(\texttt{source\_host}, \texttt{source\_log}, \texttt{line\_number}).
Without an index on \texttt{labeled\_line(source\_host, source\_log,
line\_number)}, PostgreSQL must scan the full 61{,}862-row
\texttt{labeled\_line} table for each candidate \texttt{audit\_event}.

\subsubsection{Baseline}
\begin{lstlisting}[language=SQL]
EXPLAIN (ANALYZE, BUFFERS, FORMAT TEXT)
WITH labeled_audit AS ( ... )    -- query A6 inlined
SELECT ... FROM labeled_audit ORDER BY timestamp, host_key;
\end{lstlisting}
% TODO: insert figures/explain_before.png
\textbf{Observed plan (before index).} % TBD naman: paste planning time,
% execution time, and the offending Seq Scan node from EXPLAIN ANALYZE.

\subsubsection{Index Definition}
The index lives in \texttt{sql/3nf/\_indexes.sql} and is a covering
B-tree on the natural key:

\begin{lstlisting}[language=SQL]
CREATE INDEX IF NOT EXISTS
  idx_labeled_line_source_log_line_number
ON labeled_line (source_host, source_log, line_number);
\end{lstlisting}

\subsubsection{After Index}
\begin{lstlisting}[language=SQL]
EXPLAIN (ANALYZE, BUFFERS, FORMAT TEXT)
WITH labeled_audit AS ( ... )    -- same query
SELECT ... FROM labeled_audit ORDER BY timestamp, host_key;
\end{lstlisting}
% TODO: insert figures/explain_after.png
\textbf{Observed plan (after index).} % TBD naman: planning + execution
% time, and the new Index Scan node.

\textbf{Discussion.} % TBD naman: 1-2 sentences on the speedup factor.
The textbook expectation is that the new index turns a sequential scan
into an index lookup, dropping the relevant inner-loop cost roughly
linearly with table size.

% =========================================================================
\subsection{Persisted View}
% =========================================================================

The investigation pattern from queries A4 and A6 (joining audit events,
host inventory, and attack labels into a single attack-timeline row set)
recurs across the dashboard and the report. We persist it as a
\textsc{view} so it can be queried by simple \texttt{SELECT *} from the
dashboard layer:

\begin{lstlisting}[language=SQL]
CREATE OR REPLACE VIEW v_attack_timeline AS
SELECT h.host_key,
       ae.timestamp,
       ae.type            AS audit_type,
       am.acct            AS target_account,
       am.exe             AS executable,
       am.unit            AS service_unit,
       ap.phase_name,
       string_agg(al.label_name, ', '
                  ORDER BY al.label_name)  AS labels
FROM audit_event ae
JOIN host h          ON h.host_id     = ae.host_id
JOIN labeled_line ll ON ll.source_host = h.host_key
                    AND ll.source_log  = 'audit.log'
                    AND ll.line_number = ae.line_number
JOIN labeled_line_label lll ON lll.labeled_line_id = ll.labeled_line_id
JOIN attack_label al ON al.label_id = lll.label_id
JOIN attack_phase ap ON ap.phase_id = al.phase_id
LEFT JOIN audit_message am ON am.event_id = ae.event_id
GROUP BY h.host_key, ae.timestamp, ae.type,
         am.acct, am.exe, am.unit, ap.phase_name;
\end{lstlisting}

The view is committed to \texttt{sql/3nf/\_views.sql} and is referenced
by the dashboard charts in Section~\ref{sec:dashboard}.

\subsection{Category Coverage}
Table \ref{tab:categories} maps each query to the rubric category it
demonstrates. Five of the six rubric categories are represented; the
sixth (\textsc{set operations}) was not needed for any analytical
question and we did not force a synthetic example.

\begin{table}[htbp]
\caption{Query-to-category mapping.}
\label{tab:categories}
\centering
\begin{tabular}{lll}
\toprule
\textbf{Category} & \textbf{Queries} & \textbf{Count} \\
\midrule
Subqueries           & A1                                   & 1 \\
CTEs                 & A2, A4, A5, A6, A7                   & 5 \\
Window functions     & A3, A4, A5, A7                       & 4 \\
Multi-table join (\(\geq 3\)) & B6, A4, A6, persisted view  & 4 \\
Indexing / EXPLAIN   & Section~\ref{sec:explain}            & 1 \\
Views                & \texttt{v\_attack\_timeline}         & 1 \\
\bottomrule
\end{tabular}
\end{table}
